\documentclass[11pt,a4paper]{article}
\usepackage[T1]{fontenc}
\usepackage[utf8]{inputenc}
\usepackage{lmodern}
\usepackage[a4paper,margin=1in]{geometry}
\usepackage{microtype}
\usepackage{enumitem}
\usepackage{amsmath,amssymb}
\usepackage{mathtools}
\usepackage{hyperref}

\usepackage{mathrsfs}

\hypersetup{
	colorlinks=true,
	linkcolor=black,
	urlcolor=blue,
	citecolor=black
}

\setlist[itemize]{leftmargin=1.2em}
\setlist[enumerate]{leftmargin=1.2em}

\newcommand{\Title}[1]{\section*{#1}\vspace{-0.25em}\hrule\vspace{0.8em}}
\newcommand{\Field}[2]{\textbf{#1:}~#2\par}
\newcommand{\Affil}[1]{\textit{#1}\par}
\newcommand{\Abstract}{\vspace{0.6em}\textbf{Abstract.}~}

\begin{document}
	
	\begin{center}
		{\LARGE Online Algebra, Logic, and Topology Seminar}\\[0.25em]
		{\large Category Theory Seminar — Abstract}\\[0.5em]
		January 13th, 2026
	\end{center}
	
	\vspace{1.25em}\hrule\vspace{1.25em}
	
	\Title{Ehresmann connections in a tangent category}
	
	\Field{Speaker}{Geoffrey Cruttwell}
	\Affil{\hspace{1.15cm}Mount Allison University, Canada}
	
	\noindent \Field{Date \& Time}{January 13th, 2026 (Tuesday) — 4{:}00\,PM (Europe/Lisbon time)}
	\noindent \Field{Place}{Google Meet}
	
	\Abstract
	Tangent categories, first defined by Ji\v{r}\'{\i} Rosick\'y in 1984, are a categorical structure which allows one to abstractly work with and reason about many of the ideas in differential geometry. In this talk, I'll discuss how to define and work with Ehresmann connections in a tangent category. No prior knowledge of tangent categories or differential geometry will be assumed.
	
	\vspace{1em}
	
	The content of the talk is based on joint work with Marcello Lanfranchi.
	
\end{document}
